% part: intuitionistic-logic
% chapter: introduction
% section: natural-deduction

\documentclass[../../../include/open-logic-section]{subfiles}

\begin{document}

\olfileid[zh]{int}{int}{ntd}

\olsection{自然演绎}

不含 $\FalseCl$ 规则的自然演绎是一个标准的用于直觉主义逻辑的!!{derivation}系统。我们在此复述这些规则，并用 BHK 解释来指明其动机。在
每种情形下，我们都可以把某条规则设想为允许我们断言：如果前提有构造，那么
结论也有构造。

由于自然演绎中的!!{derivation}都带有未消去的假设，我们应当把这样的!!a{derivation}，例如从!!{undischarged}假设~$\Gamma$ 推出的~$!A$，看作
把 $\Gamma$ 中所有 $!B \in \Gamma$ 的构造变为~$!A$ 的一个构造的函数。如果
存在从无!!{undischarged}假设推出的~$!A$ 的!!a{derivation}，那么在 BHK
解释的意义下就存在~$!A$ 的一个构造。不过，为了讨论方便，非必要时我们将
省去~$\Gamma$。

一个假设~$!A$ 本身就是从!!{undischarged}假设~$!A$ 推出~$!A$ 的!!a{derivation}。这符合 BHK 解释：作用于构造的恒等函数把~$!A$ 的任何构造
都变为~$!A$ 的一个构造。

\subsection{合取}

\begin{defish}
\AxiomC{$!A$}
\AxiomC{$!B$}
\RightLabel{\Intro{\land}}
\BinaryInfC{$!A \land !B$}
\DisplayProof
\hfill
\begin{tabular}{l}
\AxiomC{$!A \land !B$}
\RightLabel{\Elim{\land}}
\UnaryInfC{$!A$}
\DisplayProof
\\[3ex]
\AxiomC{$!A \land !B$}
\RightLabel{\Elim{\land}}
\UnaryInfC{$!B$}
\DisplayProof
\end{tabular}
\end{defish}

设我们分别有 $!A_1$ 与 $!A_2$ 的构造 $N_1$、$N_2$。那么我们也有 $!A_1
\land !A_2$ 的一个构造，即有序对 $\tuple{N_1, N_2}$。

在 BHK 解释下，$!A_1 \land !A_1$ 的一个构造就是一个有序对
$\tuple{N_1, N_2}$。于是设我们有这样的一个有序对。那么我们也各有每个合取支
的一个构造：$N_1$ 是~$!A_1$ 的一个构造，而 $N_2$ 是 $!A_2$ 的一个构造。

\subsection{条件句}

\begin{defish}
  \AxiomC{$\Discharge{!A}{u}$}
  \DeduceC{$!B$}
  \DischargeRule{$\Intro{\lif}$}{u}
  \UnaryInfC{$!A \lif !B$}
  \DisplayProof
\hfill
  \AxiomC{$!A \lif !B$}
  \AxiomC{$!A$}
  \RightLabel{$\Elim{\lif}$}
  \BinaryInfC{$!B$}
  \DisplayProof
\end{defish}

如果我们有从!!{undischarged}假设~$!A$ 推出~$!B$ 的!!a{derivation}，那么
就有一个函数~$f$ 把~$!A$ 的构造变为~$!B$ 的构造。这同一个函数就是 $!A
\lif !B$ 的一个构造。因此，若 $\Intro\lif$ 的前提建立在~$!A$ 的一个构造之上而具有构造，
则结论 $!A \lif !B$ 也有构造。

另一方面，设 $!A$ 的构造~$N$ 与 $!A \lif !B$ 的构造~$f$ 都存在。$!A \lif
!B$ 的一个构造就是把 $!A$ 的构造变为~$!B$ 的构造的函数。所以，$f(N)$ 就是
~$!B$ 的一个构造，即 $\Elim\lif$ 的结论有一个构造。

\subsection{析取}

\begin{defish}
\begin{tabular}{l}
\AxiomC{$!A$}
\RightLabel{\Intro{\lor}}
\UnaryInfC{$!A \lor !B$}
\DisplayProof
\\[3ex]
\AxiomC{$!B$}
\RightLabel{\Intro{\lor}}
\UnaryInfC{$!A \lor !B$}
\DisplayProof
\end{tabular}
\hfill
\AxiomC{$!A \lor !B$}
\AxiomC{$\Discharge{!A}{n}$}
\DeduceC{$!C$}
\AxiomC{$\Discharge{!B}{n}$}
\DeduceC{$!C$}
\DischargeRule{\Elim{\lor}}{n}
\TrinaryInfC{$!C$}
\DisplayProof
\end{defish}

如果我们有~$!A_i$ 的一个构造~$N_i$，就可以把它变为~$!A_1 \lor !A_2$ 的一个
构造~$\tuple{i, N_i}$。另一方面，设我们有~$!A_1 \lor !A_2$ 的一个构造，即一个
有序对 $\tuple{i, N_i}$，其中 $N_i$ 是~$!A_i$ 的一个构造，并且有函数 $f_1$、
$f_2$，它们分别把 $!A_1$、$!A_2$ 的构造变为~$!C$ 的构造。那么 $f_i(N_i)$ 就是
~$!C$ 的一个构造，即~$\Elim\lor$ 的结论。

\subsection{荒谬}

\begin{defish}
  \AxiomC{$\lfalse$}
  \RightLabel{\FalseInt}
  \UnaryInfC{$!A$}
  \DisplayProof
\end{defish}

如果我们有从!!{undischarged}假设~$!B_1$、\dots、$!B_n$
推出~$\lfalse$ 的!!a{derivation}，那么就有一个函数~$f(M_1, \dots, M_n)$ 把
$!B_1$、\dots、$!B_n$ 的构造变为~$\lfalse$ 的一个构造。由于 $\lfalse$ 没有构造，
$!B_1$、\dots、$!B_n$ 全部同时有构造也不可能。因此，$f$ 也具有如下性质：
\emph{如果} $M_1$、\dots、$M_n$ 分别是 $!B_1$、\dots、$!B_n$ 的构造，
\emph{那么} $f(M_1, \dots, M_n)$ 就是~$!A$ 的一个构造。

\subsection{$\lnot$ 的规则}

由于 $\lnot !A$ 定义为 $!A \lif \lfalse$，严格说来我们并不需要~$\lnot$ 的规则。
但如果需要，看起来就是这样的：

\begin{defish}
\AxiomC{$\Discharge{!A}{n}$}
\noLine
\DeduceC{$\lfalse$}
\DischargeRule{\Intro{\lnot}}{n}
\UnaryInfC{$\lnot !A$}
\DisplayProof
\hfill
\AxiomC{$\lnot !A$}
\AxiomC{$!A$}
\RightLabel{\Elim{\lnot}}
\BinaryInfC{$\lfalse$}
\DisplayProof
\end{defish}


\subsection{\usetoken{P}{derivation}的例子}

\begin{enumerate}
\item $\Proves !A \lif (\lnot !A \lif \lfalse)$，即 $\Proves !A
  \lif ((!A \lif \lfalse) \lif \lfalse)$
  \begin{prooftree}
    \AxiomC{$\Discharge{!A}{2}$}
    \AxiomC{$\Discharge{!A \lif \lfalse}{1}$}
    \RightLabel{\Elim\lif}
    \BinaryInfC{$\lfalse$}
    \DischargeRule{\Intro\lif}{1}
    \UnaryInfC{$(!A \lif \lfalse)\lif \lfalse$}
    \DischargeRule{\Intro\lif}{2}
    \UnaryInfC{$!A \lif (!A \lif \lfalse) \lif \lfalse$}
  \end{prooftree}

\item $\Proves ((!A \land !B) \lif !C) \lif (!A \lif (!B \lif !C))$
  \begin{prooftree}
    \AxiomC{$\Discharge{(!A \land !B) \lif !C}{3}$}
    \AxiomC{$\Discharge{!A}{2}$}
    \AxiomC{$\Discharge{!B}{1}$}
    \RightLabel{\Intro\land}
    \BinaryInfC{$!A \land !B$}
    \RightLabel{\Elim\lif}
    \BinaryInfC{$!C$}
    \DischargeRule{\Intro\lif}{1}
    \UnaryInfC{$!B \lif !C$}
    \DischargeRule{\Intro\lif}{2}
    \UnaryInfC{$!A \lif (!B \lif !C)$}
    \DischargeRule{\Intro\lif}{3}
    \UnaryInfC{$((!A \land !B) \lif !C) \lif (!A \lif (!B \lif !C))$}
  \end{prooftree}

\item $\Proves \lnot(!A \land \lnot !A)$，即 $\Proves (!A \land (!A
  \lif \lfalse))\lif \lfalse$
  \begin{prooftree}
    \AxiomC{$\Discharge{!A \land (!A \lif \lfalse)}{1}$}
    \RightLabel{\Elim\land}
    \UnaryInfC{$!A \lif \lfalse$}
    \AxiomC{$\Discharge{!A \land (!A \lif \lfalse)}{1}$}
    \RightLabel{\Elim\land}
    \UnaryInfC{$!A$}
    \RightLabel{\Elim\lif}
    \BinaryInfC{$\lfalse$}
    \DischargeRule{\Intro\lif}{1}
    \UnaryInfC{$(!A \land (!A \lif \lfalse))\lif \lfalse$}
  \end{prooftree}

\item $\Proves \lnot\lnot(!A \lor \lnot !A)$，即 $\Proves ((!A \lor
  (!A \lif \lfalse))\lif \lfalse)\lif \lfalse$
  \begin{prooftree}\hskip-3em
    \AxiomC{$\Discharge{(!A \lor (!A \lif \lfalse))\lif \lfalse}{2}$}
    \AxiomC{$\Discharge{(!A \lor (!A \lif \lfalse))\lif \lfalse}{2}$}
    \AxiomC{$\Discharge{!A}{1}$}
    \RightLabel{\Intro\lor}
    \UnaryInfC{$!A \lor (!A \lif \lfalse)$}
    \RightLabel{\Elim\lif}
    \BinaryInfC{$\lfalse$}
    \DischargeRule{\Intro\lif}{1}
    \UnaryInfC{$!A \lif \lfalse$}
    \RightLabel{\Intro\lor}
    \UnaryInfC{$!A \lor (!A \lif \lfalse)$}
    \RightLabel{\Elim\lif}
    \insertBetweenHyps{\hskip -4em}
    \BinaryInfC{$\lfalse$}
    \DischargeRule{\Intro\lif}{2}
    \UnaryInfC{$((!A \lor (!A \lif \lfalse))\lif \lfalse)\lif \lfalse$}
  \end{prooftree}
\end{enumerate}

\begin{prop}
  如果 $\Gamma \Proves !A$ 在直觉主义逻辑的自然演绎中成立，那么 $\Gamma
  \Proves !A$ 在经典逻辑中也成立。特别是，如果 $!A$ 是直觉主义定理，那么
  $!A$ 也是经典定理。
\end{prop}

\begin{proof}
  每条自然演绎规则也都是经典自然演绎中的一条规则，所以直觉主义逻辑中的每个!!{derivation}也都是经典逻辑中的!!a{derivation}。
\end{proof}

\begin{prob}
  在直觉主义逻辑中给出下列!!{formula}的!!{derivation}：
  \begin{enumerate}
    \item $(\lnot !A \lor !B) \lif (!A \lif !B)$
    \item $\lnot\lnot\lnot !A \lif \lnot !A$
    \item $\lnot\lnot (!A \land !B) \liff (\lnot\lnot !A\land \lnot \lnot !B)$
    \item $\lnot (!A \lor !B) \liff (\lnot !A \land \lnot !B)$
    \item $(\lnot !A \lor \lnot !B) \lif \lnot (!A \land !B)$
    \item $\lnot\lnot ( !A \land !B) \lif  (\lnot\lnot !A \lor
    \lnot\lnot !B)$
  \end{enumerate}
\end{prob}

\end{document}
